% -----------------------------------------------------------
\section{Adam-God}
% -----------------------------------------------------------
	\label{ADAMGOD}
	
% ----------------------
\subsection{See also}
% ----------------------

	\hyperref[CONDESCENSION]{Condescension}, \hyperref[ADAM]{Adam},

% % ----------------------
% \subsection{1828 American Dictionary of the English Language} % *1828
% % ----------------------
	% \label{ADAMGOD-1828}

	% \begin{compactitem}
		% \item
	% \end{compactitem}

% % ----------------------
% \subsection{Oxford English Dictionary} % *OED
% % ----------------------
	% \label{ADAMGOD-OED}

	% \begin{compactitem}
		% \item[]
	% \end{compactitem}

% ----------------------
\subsection{Scriptural Usage} % *SCRIPTURES
% ----------------------
	\label{ADAMGOD-SCRIPTURES}

	% % -----
	% \subsubsection{Old Testament} % *OT
	% % -----
		% \label{ADAMGOD-OT}
	
		% % --
		% \paragraph{KJV}
		% % --
			% \label{ADAMGOD-OT-KJV}
	
			% \begin{tabular}{ll}
				% Total occurrences in the OT & 
			% \end{tabular}
			
			% \begin{compactdesc}
				% \item[]
			% \end{compactdesc}
			
		% % --
		% \paragraph{Hebrew Bible}
		% % --
			% \label{ADAMGOD-HEBREW}
		
			% %
			% \subparagraph{\texorpdfstring{\he{sample word}}{}}
			% %
				% \label{PAGA} % whatever is in step bible without periods
			
				% \begin{compactdesc}
					% \item[HALOT , PDF ] \textbf{}
						% \begin{compactitem}
							% \item[1]
						% \end{compactitem}
					% \item[BDB , PDF ] \textbf{}
						% \begin{compactitem}
							% \item[1]
						% \end{compactitem}
					% \item[DCH PDF ] \textbf{}
						% \begin{compactitem}
							% \item[1]
						% \end{compactitem}
				% \end{compactdesc}
				
				% \begin{compactdesc}
					% \item[Some OT uses] \textbf{}
						% \begin{compactdesc}
							% \item[]
						% \end{compactdesc}
				% \end{compactdesc}
			
		% % --
		% \paragraph{Septuagint}
		% % --
			% \label{ADAMGOD-LXX}
		
			% %
			% \subparagraph{\texorpdfstring{\gr{sample word}}{}}
			% %
		
	% -----
	\subsubsection{New Testament} % *NT
	% -----
		\label{ADAMGOD-NT}
	
		% % --
		% \paragraph{KJV}
		% % --
			% \label{ADAMGOD-NT-KJV}
			
	
			% \begin{tabular}{ll}
				% Total occurrences in the NT & 
			% \end{tabular}
			
			\begin{compactdesc}
				\item[]
			\end{compactdesc}
			
		% --
		\paragraph{Greek New Testament} % *GREEK
		% --
			\label{ADAMGOD-GREEK}
			
			\begin{compactdesc}
				\item[{\hyperref[MATTHEW-10-24]{Matthew 10:24--25}}]
				\item[{\hyperref[LUKE-2-34]{Luke 2:34}}] \phantom{.}
					\begin{compactitem}
						\item I'm thinking specifically of the ``rising and falling'' part.
					\end{compactitem}
				\item[{\hyperref[LUKE-3-38]{Luke 3:38}}]
				\item[{\hyperref[LUKE-9-31]{Luke 9:31}}] \phantom{.}
					\begin{compactitem}
						\item See notes at the phrase ``\hyperref[ROUND-PHRASES]{one eternal round}.''
					\end{compactitem}
				\item[{\hyperref[LUKE-10-22]{Luke 10:22}}] \phantom{.}
					\begin{compactitem}
						\item Admittedly a stretch to go with Adam-God stuff.
					\end{compactitem}
				\item[{\hyperref[LUKE-15-11]{Luke 15:11--32}}] \phantom{.}
					\begin{compactitem}
						\item The prodigal son, which is so rich that it can definitely be taken along Adam-God lines, in particular with the eternal lives and multiple resurrections.1
					\end{compactitem}
				\item[{\hyperref[JOHN-3-13]{John 3:13--15}}] \gr{καὶ οὐδεὶς ἀναβέβηκεν εἰς τὸν οὐρανὸν εἰ μὴ ὁ ἐκ τοῦ οὐρανοῦ καταβάς, ὁ υἱὸς τοῦ ἀνθρώπου. 14 καὶ καθὼς Μωϋσῆς ὕψωσεν τὸν ὄφιν ἐν τῇ ἐρήμῳ, οὕτως ὑψωθῆναι δεῖ τὸν υἱὸν τοῦ ἀνθρώπου, 15 ἵνα πᾶς ὁ πιστεύων ἐν αὐτῷ ἔχῃ ζωὴν αἰώνιον.}
					\begin{compactitem}
						\item This verse may apply to Adam-God with the idea that a god must lose exaltation in order to re-enter a non-celestial world, do some activity there, and then be exalted again.
					\end{compactitem}
				\item[{\hyperref[1CORINTHIANS-15]{1 Corinthians 15}}] \phantom{.}
					\begin{compactitem}
						\item Especially the second half of this chapter; notice the connections of certain Greek words to each other, which are obscured in the English translation.
					\end{compactitem}
				\item[{\hyperref[EPHESIANS-3-14]{Ephesians 3:14--15}}] For this cause I bow my knees unto the Father of our Lord Jesus Christ, 15 Of whom the whole family in heaven and earth is named,
					\begin{compactitem}
					\item If Adam is the Father, then Paul is bowing to Adam, ``of whom the whole family in heaven and earth is named.'' Then see \hyperref[2NEPHI-9-21]{2 Nephi 9:21}: ``And he cometh into the world that he may save all men, if they will hearken unto his voice. For behold, he suffereth the pains of all men, yea, the pains of every living creature, both men women and children, which belong to \textbf{the family of Adam}''; and \hyperref[MORMOM-3-20]{Mormon 3:20}: ``And these things do the Spirit manifest unto me. Therefore I write unto you all; and for this cause I write unto you, that ye may know that ye must all stand before the judgment seat of Christ, yea, every soul which belongeth to \textbf{the whole human family of Adam} ---and ye must stand to be judged of your works, whether they be good or evil---''.
					\item There's a possible dual meaning in Eph 3:15 though. It could refer to Christ, after whom we are named when we become his sons and daugthers through the new birth; but it seems also to refer to the Father, after whom all families are named.
					\item See \hyperref[EPHESIANS-3-15]{Ephesians 3:15} for a comment on the Greek text.
					\end{compactitem}
			\end{compactdesc}
		
			% %
			% \subparagraph{\texorpdfstring{\gr{sample word}}{}}
			% %
				% \label{KATALLAGE}
			
				% \begin{compactdesc}
					% \item[BDAG , PDF ] \textbf{}
						% \begin{compactitem}
							% \item 
						% \end{compactitem}
					% \item[LSJ , PDF ] \textbf{}
						% \begin{compactitem}
							% \item 
						% \end{compactitem}
				% \end{compactdesc}
				
				% \begin{compactdesc}
					% \item[Some NT uses] \textbf{}
						% \begin{compactdesc}
							% \item[]
						% \end{compactdesc}
				% \end{compactdesc}
		
	% -----
	\subsubsection{Book of Mormon} % *BOM
	% -----
		\label{ADAMGOD-BOM}
	
		% \begin{tabular}{ll}
			% Total occurrences in the BOM & 
		% \end{tabular}
		
		\begin{compactdesc}
			\item[1 Nephi 5:11] And he beheld that they did contain the five books of Moses, which gave an account of the creation of the world and also of Adam and Eve, which was our first parents,
			\item[1 Nephi 11:18] And he said unto me: Behold, the virgin which thou seest is the mother of God after the manner of the flesh.
			\item[{\hyperref[2NEPHI-2]{2 Nephi 2}}] \textit{(Read this entire chapter with Adam-God in mind and see how you can make it work. Especially verse 25, which comes to be very beautiful.)}
			\item[{\hyperref[2NEPHI-2-25]{2 Nephi 2:25}}] Adam fell that men might be, and men are that they might have joy. \textit{(Think especially in the first half of this verse, of Adam giving up his exaltation and thereby ``falling'' in order to create material bodies for the offspring which he had already had spiritually.)}
			\item[2 Nephi 9:21] And he cometh into the world that he may save all men, if they will hearken unto his voice. For behold, he suffereth the pains of all men, yea, the pains of every living creature, both men women and children, which belong to the family of Adam.
			\item[3 Nephi 15:5] Behold, I am he that gave the law, and I am he which covenanted with my people Israel. Therefore the law in me is fulfilled, for I have come to fulfill the law. Therefore it hath an end. \textit{(I don't quite know how the Adam-godders interpret scriptures like these, where it seems like Christ is saying that he's Jehovah, or the God that the people associated with in the OT. As I understand Adam-God, there's a hierarchy going from Elohim to Jehovah and then to Michael (our Father) and Christ. So Jehovah and Christ have to be different people on this account.)}
		\end{compactdesc}
		
	% -----
	\subsubsection{Doctrine and Covenants} % *DC
	% -----
		\label{ADAMGOD-DC}
	
		% \begin{tabular}{ll}
			% Total occurrences in the DC & 
		% \end{tabular}
		
		\begin{compactdesc}
			\item[{\hyperref[DC27-11]{DC 27:11}}] 
			\item[{\hyperref[DC29-34]{DC 29:34}}] 
			\item[{\hyperref[DC78-16]{DC 78:16}}] 
			\item[{\hyperref[DC132-22]{DC 132:22--25}}]
				\begin{compactitem}
					\item I'm thinking specifically of the ``lives'' and ``deaths'' language of these verses.
				\end{compactitem}
			\item[{\hyperref[DC137-5]{DC 137:5}}] (See notes at this verse as well, it's the original source that's important; this is in the list of JS quotes below.)
		\end{compactdesc}
		
	% % -----
	% \subsubsection{Pearl of Great Price} % *PGP
	% % -----
		% \label{ADAMGOD-PGP}
	
		% \begin{tabular}{ll}
			% Total occurrences in the PGP & 
		% \end{tabular}
		
		% \begin{compactdesc}
			% \item[]
		% \end{compactdesc}
		
% % ----------------------
% \subsection{Key Phrases} % *KEYPHRASES *PHRASES
% % ----------------------
	% \label{ADAMGOD-PHRASES}

	% \begin{compactdesc}
		% \item[] \textbf{#times} | list
			% \begin{compactdesc}
				% \item[Bible anchor verses] ---
					% \begin{compactdesc}
						% \item[Romans 5:5] \gr{}
					% \end{compactdesc}
				% \item[Commentary] ---
			% \end{compactdesc}
	% \end{compactdesc}
	
% % ----------------------
% \subsection{Bible Dictionaries} % *DICTIONARIES
% % ----------------------
	% \label{ADAMGOD-BD}

% ----------------------
\subsection{Restoration Usage} % *RESTORATION
% ----------------------
	\label{ADAMGOD-RESTORATION}

	% -----
	\subsubsection{Joseph Smith} % *JS
	% -----
		\label{ADAMGOD-JS}
	
		\begin{compactdesc}
			\item[{\hyperref[EMS-1836-JS-21-JAN-1836]{Visions, 21 January 1836 [D\&C 137]}}] ...I saw the beautiful streets of that kingdom, which had the appearance of being paved with gold--- I saw father Adam, and Abraham and Michael and my father and mother, my brother Alvin that has long since slept...
				\begin{compactitem}
					\item[Commentary] After ``Michael'' in the above quote, the JSP gives the following footnote. It's a good reminder that, although this quote has clear relevance for Adam-God, we have to also think about everything else JS said as well: ``Although Adam and Michael, the archangel, are designated here as separate persons, the previous year JS approved publication of the Doctrine and Covenants, which described them as the same person: `Michael, or Adam, the father of all, the prince of all, the ancient of days.' Likewise, in a 1 January 1834 letter to John Whitmer, Oliver Cowdery wrote that he had `been informed from a proper source that the Angel Michael is no less than our father Adam.' (Doctrine and Covenants 50:2, 1835 ed. [D\&C 27:11]; Oliver Cowdery, Kirtland, OH, to John Whitmer, Missouri, 1 Jan. 1834, in Cowdery, Letterbook, 15; see also Richards, `Pocket Companion,' 74--75; and Robert B. Thompson, Sermon Notes, 5 Oct. 1840, JS Collection, CHL.)''
				\end{compactitem}
		\end{compactdesc}
		
	% -----
	\subsubsection{Temple Ordinances}
	% -----
		\label{ADAMGOD-TEMPLE}
	
		\begin{compactdesc}
			\item[{\hyperref[ENDOWMENT-alone]{The creation of Adam and Eve}}] \textit{(See \hyperref[ENDOWMENT-alone]{this note} about a possible earlier version of the endowment, and the way it describes Elohim and Jehovah.)}
			\item[{\hyperref[TEMPLE-ENDOWMENT-ENDOW-PRE-1990-GARDEN]{Adam as Lord}}] We give you dominion over all these things and make you, Adam, lord over the whole earth and all things on the face thereof. \textit{(Would this include Christ?)}
			\item[{\hyperref[TEMPLE-ENDOWMENT-ENDOW-PRE-1990-GARDEN]{Apostles of the Lord Jesus Christ}}] \textit{(After Peter, James, and John make their first visit, Elohim tells Jehovah, ``}Jehovah, instruct Peter, James, and John to go down in their true character, as apostles of the Lord Jesus Christ, to the man Adam and his posterity in the telestial world and to cast Satan out of their midst\textit{.'' Jehovah then gives them this instruction: ``}Peter, James, and John, go down in your true character, as apostles of the Lord Jesus Christ, to the man Adam and his posterity in the telestial world. Cast Satan out of their midst\textit{.'' If Jehovah is Christ, then it's a little awkward for him to speak to his apostles in this way, but it's still possible to intepret it that way.)}
			\item[{\hyperref[TEMPLE-ENDOWMENT-ENDOW-PRE-1990-LECTURE]{Newer lecture at the veil}}] The 1990 revision eliminated the lecture at the veil, which was not the same as BY's \hyperref[EMS-1877-BY-7-FEB-1877]{original lecture} but was kind of long and much more substantial than what happens now. That lecture describes what happens in the endowment in quite a bit of detail. Part of it says, speaking of the temple visitors: ``you heard the voices of persons representing a council of the Gods---Elohim, Jehovah, and Michael.'' So they are told that the voices represent a council \textit{of the Gods}---that all in the council are Gods, and Michael is one of them. Soon after the lecture also says, ``Michael, one of the council of the Gods, became the man Adam, to whom was given the woman Eve. However, as Adam he did not remember his life and labors in the council,'' again affirming that Michael belonged to the council of Gods.
		\end{compactdesc}
	
	% % -----
	% \subsubsection{Brigham Young} % *BY
	% % -----
		% \label{ADAMGOD-BY}
	
		% \begin{compactdesc}
			% \item[{\hyperref[KEY]{Date/Source}}]
		% \end{compactdesc}
		
	% -----
	\subsubsection{Lorenzo Snow} % *LS
	% -----
		\label{ADAMGOD-LS}
		
		\begin{compactdesc}
			\item[14 October 1882, WWJ 8:126] President Taylor told what Joseph Smith said to him upon that subject \& said if we do not Embrace that principle soon the Keys will be turned against us for if we do not keep the same law that Our Heavenly Father has we Cannot go with him. The word of the Lord to us was that if we did not obey that Law we Could not go whare our Heavenly Father dwells. A man obeying a lower Law is not qualified to preside over those who keep a higher Law.
		\end{compactdesc}
	
% ----------------------
\subsection{Commentary and Studies} % *COMMENTARY *STUDIES
% ----------------------
	\label{ADAMGOD-COMMENTARY}

	% -----
	\subsubsection{Key Points}
	% -----
		\label{ADAMGOD-KEYS}
	
		\begin{compactitem}
			\item See this thought on Eve, from \hyperref[EVE-KEYS]{here}: ``In the temple, Eve is called `the mother of all living' before she gives birth to a child or does anything else, really.'' One way of understanding this is that she had already been a mother, because she was already exalted.
		\end{compactitem}
		
	% -----
	\subsubsection{Thoughts on Adam-God}
	% -----
		\label{ADAMGOD-thoughts}
		
		\begin{compactenum}
			\item There is a really simple way of seeing Adam-God in the basic, canonical account of the fall---Adam and Eve were already resurrected, and had perfected, spiritual bodies, and so couldn't die. Until something changed, they were immortal, so the fall had to happen in order for them to provide mortal bodies for their offspring. They couldn't have done so before that time.
			\item Tay Reynolds once wrote the following for an assignment in the JS class in fall 2022, after doing the readings about the soul:
				\begin{compactitem}
					\item ``I am fascinated by Adam's position in this system. It seems that he is the highest angel, which also means that before the world began, he had already advanced farther than the typical spirit of a human (`Spirits can only be revealed in flaming fire, or glory. Angels have advanced farther---their light and glory being tabernacled, and hence appear in bodily shape'). If it is Adam's authority that is needed to bring the keys from heaven, wouldn't that mean that part of Christ's journey as a man required Adam's--- higher---authority? This is not meant to be incendiary, but is there a chance Adam's soul was in in a more advanced position than Christ's when this all started, at least in some ways? He is father of the human family and presides over the spirit of all men, and one of the most important parts of Christ's mission was the fact that he lived a mortal, human life, with a human spirit and body. Does that put him in Adam's stewardship? Furthermore, Christ had the keys on Earth...must He then stand before Adam in the Grand Council?
					\item ``This brings up so many questions, namely, what was Christ's soul at the beginning of all this? Was He an angel whose glory was embodied, or was Him being a less-advanced spirit a necessary condition for His ability to volunteer as Savior? How many lives/stages/reorganizations had Adam been through before He helped build the world and then became the first man? What I find confusing, in addition to everything else I've just said, is that Adam's sins as a man absolutely required Christ's redeeming power. And yet he was also this presiding archangel Michael. How can that be? It makes me think that the way position and advancement in the plan works isn't perfectly linear… like, you're a human, you're a savior, you're a god. Maybe the progression looks different for different spirits, or is simply different than we're imagining.''
				\end{compactitem}
		\end{compactenum}